\section{Einleitung}

Im deutschen Stromsystem ist der Beitrag der erneuerbaren Energien 2024 der Normalfall geworden.
Im hier verwendeten Datensatz decken sie im Jahresmittel $55{,}1\,\%$ der Last. Dieser Mittelwert
ist zugleich die am wenigsten aussagekräftige Zahl des Datensatzes, denn er entsteht aus Stunden
zwischen $12{,}0\,\%$ und $138{,}6\,\%$, aus $441$ Stunden mit negativen Börsenpreisen und aus
einzelnen Stunden über $900$\,EUR/MWh. Für das Verständnis des Gesamtsystems ist daher die zeitliche
Verteilung der Werte ausschlaggebend.

Darin liegt das Zielproblem dieses Projekts. Auf physikalischer und marktlicher Ebene beschreiben sechs kontinuierliche Zeitreihen den Zustand des Gesamtsystems, namentlich Last, Solarerzeugung, Winderzeugung, Erneuerbarenanteil, Börsenpreis sowie grenzüberschreitender Handel. Für die multidimensionale Betrachtung ganzer Tage tritt als siebtes Attribut die Anzahl negativer Preisstunden hinzu. Jede dieser Größen ist einzeln leicht verfügbar. Schwer zu beobachten ist ihr Zusammenspiel. Dieses entfaltet sich auf drei Zeitskalen gleichzeitig. Dies betrifft den Tagesgang
mit der solaren Mittagsspitze, das mehrtägige Wettergeschehen mit Windphasen und Dunkelflauten sowie
den übergeordneten Jahresgang. Etablierte Darstellungen erzwingen zumeist eine Festlegung auf eine
dieser Skalen. Während eine Jahresübersicht den Tagesgang einebnet, verliert eine isolierte Tagesansicht
den saisonalen Kontext. Die Frage \enquote{Sind negative Preise ein Solarphänomen?} erfordert jedoch beide
Skalen zugleich und zusätzlich einen Weg, ganze Tage miteinander zu vergleichen.

Daraus ergeben sich vier Fragestellungen, die mit Techniken der Informationsvisualisierung
beantwortet werden können und an denen diese Arbeit zu messen ist.

\begin{itemize}
  \item[\textbf{F1}] Wie verteilt sich die erneuerbare Deckung über Jahreszeit und Tageszeit, und
        welche der beiden Skalen erklärt mehr von der Schwankung?
  \item[\textbf{F2}] Wann treten Extremzustände auf? Gemeint sind hohe Deckung mit negativen
        Preisen auf der einen Seite und Dunkelflauten mit hohen Preisen auf der anderen. Sind das
        Einzelstunden oder Tagesmuster?
  \item[\textbf{F3}] Wie tragen Wind- und Solarenergie jeweils dazu bei? Sind sie austauschbar oder
        tragen sie zu verschiedenen Zeiten und mit verschiedenen Folgen bei?
  \item[\textbf{F4}] Wie hängen Erneuerbarenanteil, Preis und Nettohandel an einem konkreten Tag
        zusammen, und wie typisch ist dieser Tag im Vergleich zu allen anderen?
\end{itemize}

Keine dieser Fragen lässt sich mit einer einzelnen Kennzahl oder einem isolierten Diagramm
beantworten. So erfordert F1 zwei Zeitskalen gleichzeitig, F2 das Auffinden seltener Ereignisse in
$8\,690$ Messwerten und F3 sowie F4 den Vergleich mehrerer Attribute mit unterschiedlichen physikalischen
Einheiten. Ein reales Anwendungsproblem, für das Visualisierungstechniken ausgewählt, kombiniert und
begründet werden, entspricht dem Konzept einer \emph{Design Study} nach Munzner \cite{Munzner2008}.

\subsection{Anwendungshintergrund}

Für die Einordnung der Fragestellungen sind vier Zusammenhänge des Anwendungsfeldes maßgeblich.

\textbf{Dargebotsabhängige Erzeugung.} Solar- und Windanlagen speisen abhängig vom Wettergeschehen
ein, losgelöst vom aktuellen Strombedarf. Die Photovoltaik folgt einem deterministischen Tages- und
Jahresgang mit dem Maximum um die Mittagszeit und deutlichem Sommerüberschuss. Die Windenergie ist
wetterbestimmt, variiert über mehrtägige Wetterphasen und verzeichnet ein Winterhoch. Beide Profile
ergänzen einander komplementär, was für die Auswahl der Visualisierungen richtungsweisend ist.

\textbf{Residuallast und Dunkelflaute.} Der nach Abzug der erneuerbaren Einspeisung verbleibende Teil
der Last heißt Residuallast und muss von steuerbaren Kraftwerken, Speichern oder Importen gedeckt
werden. Fallen Wind und Sonne über Stunden bis Tage gleichzeitig aus, während die Last hoch ist
(\enquote{Dunkelflaute}), wird das gesamte konventionelle Reservoir gebraucht. Solche Lagen sind
selten und für die Systemauslegung dennoch prägend.

\textbf{Merit-Order und negative Preise.} Am Day-Ahead-Markt der Gebotszone DE-LU bestimmt das
Grenzkraftwerk den einheitlichen Börsenpreis. Erneuerbare Anlagen speisen mit Grenzkosten nahe null
ein und verschieben die Angebotskurve nach rechts. Ein hohes erneuerbares Angebot senkt entsprechend
den Preis. Übersteigt das Angebot die Nachfrage bei unflexibler Abregelung, fallen die Preise unter
null. Negative Börsenpreise bilden damit das reguläre Marktsignal für einen akuten Erzeugungsüberschuss.

\textbf{Grenzüberschreitender Handel.} Überschüsse fließen in das europäische Verbundnetz ab,
während Defizite durch Importe ausgeglichen werden. Der Nettohandelswert zeigt als dritter Baustein
neben Erzeugung und Preis, ob ein Überschuss im Inland verbleibt oder abfließt.

In der öffentlichen Debatte treten diese Zusammenhänge meist als Behauptungen auf
(\enquote{Erneuerbare machen Strom billig}, \enquote{im Winter liefern sie nichts}), werden aber
selten an Daten geprüft. Genau diese Prüfung soll die Anwendung ermöglichen.

\subsection{Zielgruppen}

\textbf{Z1 (Studierende und Lehrende energienaher Fächer, primär).} Als primäre Zielgruppe bearbeiten
sie praxisnahe Aufgabenstellungen und suchen konkrete Beispieltage für Seminare oder Abschlussarbeiten.
Fachbegriffe wie Merit-Order, Residuallast sowie die Einheiten GW und EUR/MWh sind ihnen geläufig,
der Umgang mit Rohdaten oder relationalen Datenbanken hingegen selten. Ihr Informationsbedarf besteht in
reproduzierbaren Fallbeispielen mit exakten Kennwerten.

\textbf{Z2 (Fachjournalismus, Energieagenturen und kommunale Versorger).} Diese Gruppe prüft und bewertet
energiewirtschaftliche Aussagen, häufig unter Zeitdruck und für eine breite Öffentlichkeit. Sie beherrscht
den Umgang mit statistischen Kennzahlen, besitzt jedoch wenig Erfahrung in komplexer Zeitreihenanalyse.
Ihr Ziel liegt im schnellen Auffinden belastbarer Einzelfälle samt Kontext, etwa dem teuersten Tag des
Jahres am 12.\ Dezember mit durchschnittlich $395{,}7$\,EUR/MWh.

\textbf{Z3 (Fachlich interessierte Prosumer).} Betreiber privater Photovoltaikanlagen, Batteriespeicher
oder Wärmepumpen mit dynamischen Stromtarifen kennen den eigenen Erzeugungsverlauf aus Monitoring-Apps.
Ihr Interesse gilt dem Abgleich dieser persönlichen Beobachtungen mit dem makroökonomischen Gesamtsystem.

Die drei Gruppen bringen etablierte Sehgewohnheiten mit, an die die Anwendung gezielt anschließt. Von
\emph{Energy-Charts} \cite{EnergyCharts} und Agora ist die Konvention \enquote{Erneuerbare in Grün- und Blautönen}
vertraut. Die Anwendung übernimmt Grün und Blau für die Erneuerbaren-Skala und reserviert ein warmes Rot für
den Preis, damit beide Größen nicht verwechselt werden. Kalender- und Aktivitätsheatmaps sind als
\enquote{eine Zelle = ein Zeitabschnitt, dunkler = mehr} allgemein bekannt. Darauf baut die Stunden-Heatmap auf.
Parallele Koordinaten setzen wir dagegen nicht voraus. Sie sind die einzige Technik, die erklärt werden muss,
weshalb über der Ansicht ein einzeiliger Hinweis steht und der Erstkontakt über Hover mit Tooltip erfolgt.
Preise werden in EUR/MWh erwartet und nicht in ct/kWh, Zeitangaben in UTC wie in den Quelldaten. Alle Zahlen
erscheinen mit deutschem Dezimalkomma.

\subsection{Überblick und Beiträge}

Die Grundlage der Anwendung bilden die EnergyCharts-Daten \cite{EnergyCharts} für Deutschland 2024,
abgerufen über den PostgREST-Spiegel der Universität Halle. Aus zwei Quelltabellen für Erzeugung und Preis
entsteht ein Datensatz mit $8\,690$ Stundenwerten und $364$ Tagesprofilen. Drei Techniken aus den vorgegebenen
Bereichen sind eng mit den Analyseaufgaben verzahnt. Eine Zeitreihen-Übersicht über die zwölf Monate dient als
Einstieg und Filter (Bereich Scatterplot und Zeitreihen), eine pixelorientierte Stunden-Heatmap fungiert als
Tag-Stunden-Matrix für alle $8\,690$ Stundenwerte (Bereich Pixel- und Icon-Techniken), und parallele
Koordinaten über sieben Tagesattribute decken den Bereich mehrdimensionaler Darstellungen ab. Die drei
Ansichten sind über einen gemeinsamen Interaktionszustand gekoppelt.

Diese Arbeit leistet fünf wesentliche Beiträge.

\begin{itemize}
  \item[\textbf{B1}] Eine reproduzierbare Datenaufbereitung (Kapitel~\ref{sec:daten}), die
        das in der Quelle unvollständige Lastfeld rechnerisch rekonstruiert und die $94$ fehlenden
        Stunden transparent als Datenlücke darstellt.
  \item[\textbf{B2}] Eine aus den Anwendungsaufgaben begründete Technikauswahl
        (Abschnitt~\ref{sec:praesentation}), bei der jede Technik gegen drei echte Alternativen
        abgewogen wird, die dasselbe mentale Modell aufbauen würden.
  \item[\textbf{B3}] Eine durchgängige Kopplung der drei Ansichten
        (Abschnitt~\ref{sec:interaktion}) über Mehrfach-Monatsfilter, geteilten Hover-Zustand,
        Mehrfachauswahl von Tagen und UND-verknüpftes Achsen-Brushing.
  \item[\textbf{B4}] Drei belegte Anwendungsfälle (Kapitel~\ref{sec:anwendungsfaelle}),
        darunter der Nachweis, dass die Häufigkeit negativer Börsenpreise vor allem dem
        Solaranteil folgt und $69{,}6\,\%$ dieser Stunden in das Fenster von 10 bis 15\,Uhr fallen.
  \item[\textbf{B5}] Eine offene Diskussion der Grenzen bei Lastrekonstruktion,
        Vorzeichenkonvention des Handelswertes und Datenlücken, damit die gezeigten Muster nicht
        überinterpretiert werden.
\end{itemize}

\abb{uebersicht}{1.0}{Gesamtansicht der Anwendung: Monatsfilter und Kennzahlen oben, darunter die
drei gekoppelten Visualisierungen, rechts das Detailpanel.}
{Screenshot der vollständigen Anwendung im Startzustand (kein Monatsfilter): Kopfzeile mit
Leitfrage, Monatsschaltflächen, sechs Kennzahlenkarten, Zeitreihen-Übersicht, Stunden-Heatmap,
parallele Koordinaten und rechts das Tagesdetailpanel.}
